\documentclass[10pt,a4paper]{amsart}
\usepackage[margin=19mm]{geometry}
\usepackage{amsmath,amssymb,mathtools}
\usepackage[scr=rsfs,cal=euler]{mathalfa}
\usepackage{xfrac}
\usepackage[unicode,hidelinks,bookmarks=false]{hyperref}
\newcommand{\ie}{i.e.\ }
\newcommand{\cf}{cf.\ }
\newcommand{\resp}{respectively\ }
\newcommand{\Subsection}{\subsection}
\providecommand{\nobreakdash}{\nobreak\hbox{-}}
\setlength{\emergencystretch}{2em}
\multlinegap0pt
\usepackage{bm}
\usepackage{bbm}
\usepackage{dsfont}
\usepackage{xspace}
\usepackage{cancel}
\usepackage{mathtools}
\usepackage{amsthm}
\makeatletter
\@ifundefined{thm}{\newtheorem{thm}{Thm}}{}
\@ifundefined{definition}{\newtheorem{definition}[thm]{Definition}}{}
\@ifundefined{proposition}{\newtheorem{proposition}[thm]{Proposition}}{}
\makeatother
\title[Voltage quantum graphs and a Gross-Tucker theorem for quantu]{Voltage quantum graphs and a Gross-Tucker theorem for quantum graphs}
\author{Bj\"orn Sch\"afer, Mariusz Tobolski}
\date{Source: arXiv:2602.20996v1 (CC BY 4.0)}
\begin{document}
\maketitle

\noindent\emph{Source and licence.} Voltage quantum graphs and a Gross-Tucker theorem for quantum graphs. Version arXiv:2602.20996v1 (CC BY 4.0), distributed under the Creative Commons Attribution 4.0 International licence, as recorded in the arXiv licence metadata for this exact version and shown on the abstract page https://arxiv.org/abs/2602.20996v1. Full rights notice: licenses/arxiv-2602.20996-CC-BY-4.0.txt.\par\smallskip

\noindent\textit{Editorial scope.} Counted unit: the Proposition whose source label is \texttt{land::prop:quantum\_isomorphism\_between\_crossed\_product\_and\_tensor\_product} in the article (source0.tex lines 1309--1324, source label \texttt{land::prop:quantum\_isomorphism\_between\_crossed\_product\_and\_tensor\_product}), delivered together with the article's own complete original proof of it (source0.tex lines 1326--1416, one \texttt{proof} environment of 91 lines). No mathematics is written, repaired or re-derived: statement and proof are the article's own text line for line. Required context retained uncounted: pre::def:qisomoprhisms\_of\_qsets\_and\_qgraphs (lines 320--341); land::eq:formula\_for\_scalar\_product\_on\_crossed\_product\_algebra (lines 495--504). Every definition, display and label of those lines is retained unchanged. Omitted from this package and not redistributed: 1--319, 342--494, 505--1308, 1325--1325, 1417--1946. Nothing inside a retained excerpt is omitted or altered: every retained body is delivered exactly as the article's lines are written, so no omission receipt of any kind accompanies this package. Citations: the retained excerpts carry no citation command at all, so no bibliography entry of the article is transcribed and no reference is redistributed. Reference adaptation: none was needed and none was made; every reference inside the delivered unit resolves inside it, so no unresolved reference or citation is produced. Printed slips kept literal: no printed word, symbol, hypothesis or number of the counted statement or of its proof was altered. Layout and class: the article's own class is \texttt{amsart} with packages including hyperref, cleveref, amsmath, amssymb, enumitem, fontenc, tikz, string-diagrams, biblatex; the wrapper is the lane's pinned \texttt{amsart} editorial wrapper at 10pt on A4, which restates the theorem-like environments the retained text needs and adds the article's own macro definitions that the retained text needs (none), with extra wrapper packages (bm, bbm, dsfont, xspace, cancel, mathtools). The delivered numbering is the unit's own. Assets: no retained span contains a figure environment, a graphics-inclusion command, a PDF-page inclusion, a table or a TikZ picture; no image file is copied into this package, the package declares no asset dependency and no image file is redistributed with it. Licence: version arXiv:2602.20996v1 of this article is distributed under the Creative Commons Attribution 4.0 International licence, as recorded in the arXiv licence metadata for this exact version and shown on the abstract page https://arxiv.org/abs/2602.20996v1. A full rights notice is delivered at licenses/arxiv-2602.20996-CC-BY-4.0.txt. Characters: every retained excerpt is pure ASCII, so the delivered engine, XeLaTeX with its default Latin Modern text font, reports no missing character.
\begin{definition}
    \label{pre::def:qisomoprhisms_of_qsets_and_qgraphs}
    Let $(B_1, \psi_1)$ and $(B_2, \psi_2)$ be quantum sets and let $A_1: B_1 \to B_1$ and $A_2: B_2 \to B_2$ be quantum graphs on $(B_1, \psi_1)$ and $(B_2, \psi_2)$, respectively.
    \begin{enumerate}
        \item The two quantum sets are \emph{isomorphic} if there is a $*$-isomorphism $\varphi: B_1 \to B_2$ such that $\psi_2 \varphi = \psi_1$. We denote by $\mathrm{Aut}(B, \psi)$ the group of quantum set automorphisms of a quantum set $(B, \psi)$.
        \item The two quantum graphs are \emph{isomorphic} if there is an isomorphism $\varphi: B_1 \to B_2$ of the underlying quantum sets such that $\varphi A_1 = A_2 \varphi$. We denote by $\mathrm{Aut}(A)$ the group of quantum graph automorphisms of a quantum graph $A$.
        \item The two quantum sets are \emph{quantum isomorphic}, written $(B_1, \psi_1) \cong_q (B_2, \psi_2)$, if there is a finite-dimensional Hilbert space $\mathcal{H}$ and a unital $\ast$-homomorphism  
        \begin{align*}
            \rho: B_1 \to B_2 \otimes B(\mathcal{H}),
        \end{align*}
        such that the map
        \begin{align*}
            p: L^2(B_1) \otimes \mathcal{H} \to L^2(B_2) \otimes \mathcal{H}, \quad a \otimes \xi \mapsto \rho(a)(1 \otimes \xi)
        \end{align*}
        is unitary as an operator between Hilbert spaces.
        \item The two quantum graphs are \emph{quantum isomorphic}, written $A_1 \cong_q A_2$, if there is a quantum isomorphism $\rho: B_1 \to B_2 \otimes B(\mathcal{H})$ between the underlying quantum sets such that
        \begin{align*}
            \rho A_1 = (A_2 \otimes \mathrm{Id}_{B(\mathcal{H})}) \rho
        \end{align*}
        holds.
    \end{enumerate}
\end{definition}

    \begin{align}
        \label{land::eq:formula_for_scalar_product_on_crossed_product_algebra}
        \left\langle b u_\chi, c u_\xi \right\rangle_{\psi_\rtimes} 
        %&= n \delta_{\chi=\xi} \tilde{\psi}(\hat{\alpha}_{\chi^{-1}}(b^\ast c))
        &= n \delta_{\chi=\xi} \langle b, c \rangle_{\tilde{\psi}}, \\
        \label{land::eq:comultiplication_formula_on_crossed_product}
        m_\rtimes^\ast\left(b u_\chi\right) &= \frac{1}{n} \sum_{\xi \in \hat{G}} \sum_i b_i^{(1)} u_\xi \otimes \hat{\alpha}_{\xi^{-1}}(b_i^{(2)}) u_{\xi^{-1} \chi}, \\
        \label{eq:formula_for_E*}
        E_\rtimes^\ast(b) &= b,
    \end{align}

\begin{proposition}
    \label{land::prop:quantum_isomorphism_between_crossed_product_and_tensor_product}
    Let $\hat{\alpha}: \hat{G} \to \mathrm{Aut}(\tilde{B}, \tilde{\psi})$ be an action of the dual of a finite abelian group $\hat{G}$ on a quantum set $(\tilde{B}, \tilde{\psi})$, and let $\mathcal{H} := L^2(G)$.
    Then, the map
    \begin{align*}
        \rho: \left\{ \;
        \begin{aligned}   
            \tilde{B} \otimes C(G) &\to \left(\tilde{B} \rtimes_{\hat{\alpha}} \hat{G}\right) \otimes B(\mathcal{H}), \\
            b \otimes \tilde{u}_\chi &\mapsto \sum_{\zeta \in \hat{G}} u_\zeta b u_\zeta^\ast u_\chi \otimes E_{\zeta, \chi^{-1} \zeta}
        \end{aligned}\right.
    \end{align*}
    is a quantum isomorphism
    \begin{align*}
        (\tilde{B}, \tilde{\psi}) \rtimes_{\hat{\alpha}} \hat{G} \cong_q (\tilde{B}, \tilde{\psi}) \otimes (C(G), \psi_G).
    \end{align*}
\end{proposition}

\begin{proof}
    To prove that $\rho$ is a $\ast$-homomorphism, we construct it via the universal property of the tensor product.
    Let $(\hat{u}_\chi)_{\chi \in \hat{G}}$ be the natural representation of $\hat{G}$ on $\mathcal{H}$ given by 
    \begin{align*}
        \hat{u}_\chi e_\xi = e_{\chi \xi} \quad \text{ for all } \chi, \xi \in \hat{G},
    \end{align*}
    and consider the maps
    \begin{align*}
        \rho_1: \tilde{B} &\to \left(\tilde{B} \rtimes_{\hat{\alpha}} \hat{G}\right) \otimes B(\mathcal{H}), & b &\mapsto \sum_{\chi \in \hat{G}} u_\chi b u_\chi^\ast \otimes E_{\chi, \chi}, \\
        \rho_2: C(G) &\to \left(\tilde{B} \rtimes_{\hat{\alpha}} \hat{G}\right) \otimes B(\mathcal{H}), & \tilde{u}_\chi &\mapsto u_\chi \otimes \hat{u}_\chi.
    \end{align*}
    It is not hard to see that $\rho_1$ and $\rho_2$ are unital $\ast$-homomorphisms. Furthermore, they have commuting ranges since for all $b \in \tilde{B}$ and $\chi \in \hat{G}$,
    \begin{align*}
        \rho_2(\tilde{u}_\chi) \rho_1(b) \rho_2(\tilde{u}_\chi)^\ast 
            &= \sum_{\xi \in \hat{G}} u_\chi u_\xi b u_\xi^\ast u_\chi^\ast \otimes \hat{u}_\chi E_{\xi, \xi} \hat{u}_\chi^\ast \\
            &= \sum_{\zeta \in \hat{G}} u_\zeta b u_\zeta^\ast \otimes E_{\zeta, \zeta}
            = \rho_1(b).
    \end{align*}
    Therefore, the universal property of the tensor product of C*-algebras yields a unique unital $\ast$-homomorphism
    \begin{align*}
        \rho: \tilde{B} \otimes C(G) \to \left(\tilde{B} \rtimes_{\hat{\alpha}} \hat{G}\right) \otimes B(\mathcal{H}),
    \end{align*}
    such that $\rho|_{\tilde{B}} = \rho_1$ and $\rho|_{C(G)} = \rho_2$. A short calculation confirms that this $\ast$-homomorphism is the map $\rho$ from the statement.

    It remains to show that the associated map 
    \begin{align*}
        p: L^2(\tilde{B} \otimes C(G)) \otimes \mathcal{H} \to L^2(\tilde{B} \rtimes_{\hat{\alpha}} \hat{G}) \otimes \mathcal{H}
    \end{align*}
    from Definition \ref{pre::def:qisomoprhisms_of_qsets_and_qgraphs} is unitary. First, let us show that the adjoint of $p$ is given by
    \begin{align}
        \label{qiso::eq:formula_for_p_adjoint}
        p^\ast: \left\{
        \begin{aligned} 
            L^2(\tilde{B} \rtimes_{\hat{\alpha}} \hat{G}) \otimes \mathcal{H} &\to L^2(\tilde{B} \otimes C(G)) \otimes \mathcal{H}, \\
            b_\chi u_\chi \otimes e_\xi &\mapsto u_\xi^\ast b_\chi u_\xi \otimes \tilde{u}_\chi \otimes e_{\chi^{-1}\xi}.
        \end{aligned} \right.
    \end{align}
    Indeed, for all $c \in \tilde{B}$ and $\alpha, \beta \in \hat{G}$ we have
    \begin{align*}
        p(c \otimes \tilde{u}_\alpha \otimes e_\beta) 
        &= \sum_{\gamma \in \hat{G}} u_\gamma c u_\gamma^\ast u_\alpha \otimes E_{\gamma, \alpha^{-1} \gamma} e_\beta
        = u_{\alpha\beta} c u_{\beta}^\ast \otimes e_{\alpha\beta}.
    \end{align*}
    It follows for all $b_\chi \in \tilde{B}$ and $\chi, \xi \in \hat{G}$,
    \begin{align*}
        \begin{aligned}
        &\hspace{-.5cm} \langle b_\chi u_\chi \otimes e_\xi, p(c \otimes \tilde{u}_\alpha \otimes e_\beta) \rangle_{\psi_{\rtimes} \otimes \psi_G} \\
            &= \langle b_\chi u_\chi \otimes e_\xi,  u_{\alpha\beta } c u_{\beta}^\ast \otimes e_{\alpha\beta} \rangle_{\psi_{\rtimes} \otimes \psi_G} \\
            &= n \delta_{\xi=\alpha\beta} \langle b_\chi u_\chi, u_{\alpha\beta } c u_{\beta}^\ast \rangle_{\psi_\rtimes}.
    \end{aligned}
    \end{align*}
    Using formula \eqref{land::eq:formula_for_scalar_product_on_crossed_product_algebra} for $\langle \cdot, \cdot \rangle_{\psi_\rtimes}$, the previous expression becomes
    \begin{align}
        \label{land::eq:computation_of_p_to_compare_with_p_adjoint}
        \begin{aligned}
            \langle b_\chi u_\chi \otimes e_\xi, p(c \otimes \tilde{u}_\alpha \otimes e_\beta) \rangle_{\psi_{\rtimes} \otimes \psi_G}
            &= n^2 \delta_{\xi=\alpha\beta} \delta_{\chi=\alpha} \langle b_\chi, \hat{\alpha}_{\alpha\beta}(c) \rangle_{\tilde{\psi}}.
        \end{aligned}
    \end{align}
    On the other hand, for the map $p^\ast$ from \eqref{qiso::eq:formula_for_p_adjoint} we compute
    \begin{align*}
        \begin{aligned}
        \langle p^\ast(b_\chi u_\chi \otimes e_\xi), &c \otimes \tilde{u}_\alpha \otimes e_\beta \rangle_{\tilde{\psi} \otimes \psi_G \otimes \psi_G} \\
            &= \langle u_\xi^\ast b_\chi u_\xi \otimes \tilde{u}_\chi \otimes e_{\chi^{-1}\xi}, c \otimes \tilde{u}_\alpha \otimes e_\beta \rangle_{\tilde{\psi} \otimes \psi_G \otimes \psi_G} \\
            &= \langle u_\xi^\ast b_\chi u_\xi, c \rangle_{\tilde{\psi}} \langle \tilde{u}_\chi, \tilde{u}_\alpha \rangle_{\psi_G} \langle e_{\chi^{-1}\xi}, e_\beta \rangle_{\psi_G} \\
            &= n^2 \delta_{\chi=\alpha} \delta_{\chi^{-1} \xi = \beta} \langle u_\xi^\ast b_\chi u_\xi, c \rangle_{\tilde{\psi}}.
        \end{aligned}
    \end{align*}
    Using $\tilde{\psi} \circ \hat{\alpha}_\xi = \tilde{\psi}$ it is not hard to see
    \begin{align*}
        \langle u_\xi^\ast b_\chi u_\xi, c \rangle_{\tilde{\psi}} 
            = \tilde{\psi}(u_\xi^\ast b_\chi^\ast u_\xi c)
            = \tilde{\psi}(b_\chi^\ast u_\xi c u_\xi^\ast)
            = \langle b_\chi, \hat{\alpha}_\xi(c) \rangle_{\tilde{\psi}}.
    \end{align*}
    Since $\chi^{-1} \xi = \beta$ and $\chi=\alpha$ entail $\xi = \alpha \beta$ we arrive at
    \begin{align}
        \langle p^\ast(b_\chi u_\chi \otimes e_\xi), c \otimes \tilde{u}_\alpha \otimes e_\beta \rangle_{\tilde{\psi} \otimes \psi_G \otimes \psi_G}
            = n^2 \delta_{\chi=\alpha} \delta_{\xi = \alpha \beta} \langle b_\chi, \hat{\alpha}_{\alpha \beta}(c) \rangle_{\tilde{\psi}}.
    \end{align}
    This is the same expression as in \eqref{land::eq:computation_of_p_to_compare_with_p_adjoint}; consequently $p^\ast$ is indeed the adjoint of the map $p$.

    Finally, we are ready to show that $p$ is a unitary map. It suffices to prove $p p^\ast = \mathrm{Id}$. For this, let $b_\chi u_\chi \in \tilde{B} \rtimes_{\hat{\alpha}} \hat{G}$ and $\xi \in \hat{G}$ and observe
    \begin{align*}
        pp^\ast(b_\chi u_\chi \otimes e_\xi) 
            &= p\left( u_\xi^\ast b_\chi u_\xi \otimes \tilde{u}_\chi \otimes e_{\chi^{-1}\xi} \right) \\
            &=  \sum_{\gamma \in \hat{G}} u_\gamma u_\xi^\ast b_\chi u_\xi u_\gamma^\ast u_\chi \otimes E_{\gamma, \chi^{-1} \gamma} e_{\chi^{-1}\xi} \\
            &= b_\chi u_\chi \otimes e_\xi.
    \end{align*}
    This finishes the proof.
\end{proof}

\end{document}
